% Section: Truy hồi nguồn (Source Retrieval)
% Yêu cầu: \usepackage{booktabs}
% Dùng: % Section: Truy hồi nguồn (Source Retrieval)
% Yêu cầu: \usepackage{booktabs}
% Dùng: % Section: Truy hồi nguồn (Source Retrieval)
% Yêu cầu: \usepackage{booktabs}
% Dùng: % Section: Truy hồi nguồn (Source Retrieval)
% Yêu cầu: \usepackage{booktabs}
% Dùng: \input{report_retrieval_section}

\section{Truy hồi nguồn (Source Retrieval)}

\subsection{Bài toán và phương pháp}

Với mỗi tài liệu nghi vấn, việc đối chiếu chi tiết trực tiếp với toàn bộ $7\,949$
tài liệu nguồn là bất khả thi về chi phí tính toán. Tầng truy hồi có nhiệm vụ thu hẹp
kho nguồn về một tập nhỏ các ứng viên khả nghi: đầu vào là tài liệu nghi vấn cùng chỉ mục
của kho, đầu ra là danh sách $k$ nguồn được xếp hạng theo mức tương đồng.

Phương pháp được sử dụng là TF-IDF kết hợp độ đo cosine. Mỗi tài liệu được biểu diễn bằng
một vector TF-IDF trên $20\,000$ ký tự đầu (cắt theo ký tự, mỗi tài liệu một vector), với
cấu hình \texttt{max\_features} $=100\,000$, \texttt{sublinear\_tf}, và loại bỏ stop-word
tiếng Anh. Vector của tài liệu nghi vấn được so cosine với toàn bộ ma trận nguồn để lấy
$k$ tài liệu có điểm cao nhất.

\subsection{Ví dụ minh hoạ biểu diễn TF-IDF}

Để làm rõ cách TF-IDF đánh trọng số cho từ, xét một kho đồ chơi gồm bốn văn bản ngắn:

\begin{itemize}
  \item[$d_1$] (truy vấn): \emph{deep neural networks classify images with high accuracy}
  \item[$d_2$:] \emph{neural networks are trained on large datasets}
  \item[$d_3$:] \emph{the model reduces the computational cost of training}
  \item[$d_4$:] \emph{convolutional networks improve image classification accuracy}
\end{itemize}

Sau khi loại stop-word tiếng Anh, mỗi từ được gán trọng số theo công thức của thư viện
\texttt{scikit-learn}:
\[
\mathrm{idf}(t) = \ln\frac{1+N}{1+\mathrm{df}(t)} + 1,
\qquad
w(t,d) = \big(1 + \ln \mathrm{tf}(t,d)\big)\cdot \mathrm{idf}(t),
\]
trong đó $N=4$ là số văn bản, $\mathrm{df}(t)$ là số văn bản chứa từ $t$, và $\mathrm{tf}(t,d)$ là
tần suất từ $t$ trong văn bản $d$; vector cuối được chuẩn hoá $L_2$. Bảng~\ref{tab:tfidf-demo}
liệt kê trọng số TF-IDF của câu truy vấn $d_1$.

\begin{table}[htbp]
\centering
\renewcommand{\arraystretch}{1.2}
\begin{tabular}{lccc}
\toprule
\textbf{Từ} & \textbf{df} & \textbf{idf} & \textbf{TF-IDF ($d_1$)} \\
\midrule
\emph{classify} & 1 & 1.916 & \textbf{0.421} \\
\emph{deep}     & 1 & 1.916 & \textbf{0.421} \\
\emph{images}   & 1 & 1.916 & \textbf{0.421} \\
\emph{high}     & 1 & 1.916 & \textbf{0.421} \\
\emph{accuracy} & 2 & 1.511 & 0.332 \\
\emph{neural}   & 2 & 1.511 & 0.332 \\
\emph{networks} & 3 & 1.223 & 0.269 \\
\bottomrule
\end{tabular}
\caption{Trọng số TF-IDF của câu truy vấn $d_1$ trên kho đồ chơi bốn văn bản ($N=4$),
sắp theo trọng số giảm dần.}
\label{tab:tfidf-demo}
\end{table}

Vì mỗi từ chỉ xuất hiện một lần trong $d_1$, thành phần $\mathrm{tf}$ là đồng đều, nên chênh
lệch trọng số hoàn toàn do IDF quyết định. Các từ đặc trưng chỉ có mặt ở một văn bản
(\emph{classify}, \emph{deep}, \emph{images}, \emph{high}; $\mathrm{df}=1$) nhận trọng số cao
nhất ($0.421$), trong khi từ phổ biến \emph{networks} ($\mathrm{df}=3$) chỉ còn $0.269$. Nhờ
cơ chế này, độ đo cosine giữa hai tài liệu bị chi phối bởi các \emph{thuật ngữ đặc trưng dùng
chung}; đó chính là tín hiệu giúp tầng truy hồi xác định đúng nguồn.

\subsection{Lựa chọn TF-IDF và kết quả so sánh}

TF-IDF được lựa chọn sau khi so sánh trực tiếp với phương pháp biểu diễn ngữ nghĩa (neural
embedding) trên tập validation gồm $5\,522$ tài liệu nghi vấn. Kết quả được trình bày trong
Bảng~\ref{tab:retrieval-recall}: TF-IDF vượt trội ở mọi mức Recall@k cũng như MRR.

\begin{table}[htbp]
\centering
\renewcommand{\arraystretch}{1.2}
\begin{tabular}{lcccc}
\toprule
\textbf{Phương pháp} & \textbf{R@1} & \textbf{R@5} & \textbf{R@10} & \textbf{MRR} \\
\midrule
\textbf{TF-IDF (từ vựng)}       & \textbf{0.975} & \textbf{0.994} & \textbf{0.995} & \textbf{0.985} \\
Neural (bge-small)              & 0.908          & 0.965          & 0.975          & 0.934          \\
\bottomrule
\end{tabular}
\caption{Recall@k và MRR trên tập validation ($5\,522$ tài liệu nghi vấn).}
\label{tab:retrieval-recall}
\end{table}

Sự vượt trội của TF-IDF xuất phát từ hai yếu tố. Thứ nhất, đạo văn trong bộ PAN có độ trùng
lặp từ vựng cao, và các tài liệu khoa học chia sẻ nhiều thuật ngữ chuyên ngành, tên riêng,
ký hiệu và con số; đây là các tín hiệu mà mô hình túi từ (bag-of-words) nắm bắt trực tiếp.
Thứ hai, biểu diễn ngữ nghĩa \texttt{bge-small} bị chi phối bởi chủ đề: các tài liệu cùng
lĩnh vực nằm gần nhau trong không gian embedding, làm mờ ranh giới giữa nguồn thật và các
tài liệu chỉ trùng chủ đề, dẫn đến hạng của nguồn đúng bị hạ thấp.

\subsection{Giới hạn của biểu diễn từ vựng}

Do là mô hình túi từ, TF-IDF chỉ nắm bắt phần từ vựng bề mặt được dùng chung giữa hai văn
bản. Cần phân biệt hai dạng che giấu (obfuscation) với ảnh hưởng khác nhau:

\begin{itemize}
  \item \textbf{Đảo trật tự câu hoặc mệnh đề} không ảnh hưởng đến TF-IDF, vì biểu diễn túi từ
  bất biến với thứ tự từ: toàn bộ từ được giữ nguyên nên độ trùng từ vựng không đổi (ví dụ~4
  trong Bảng~\ref{tab:paraphrase}, Jaccard $=1.00$).
  \item \textbf{Thay từ đồng nghĩa, dịch đi--dịch lại (back-translation) và viết lại ý} là giới
  hạn thực sự: khi các từ dùng chung bị thay bằng từ khác về mặt hình thức, độ trùng từ vựng
  giảm dần về $0$ và tín hiệu lexical biến mất. Ở các đoạn bị che giấu theo cách này, TF-IDF
  không còn khả năng phát hiện.
\end{itemize}

Bảng~\ref{tab:paraphrase} minh hoạ mức suy giảm độ trùng từ vựng qua các kỹ thuật đạo văn,
đo bằng hai chỉ số: Jaccard trên tập từ nội dung, và cosine giữa hai vector TF-IDF (cùng cấu
hình với tầng truy hồi). Khi từ vựng dùng chung bị thay thế dần (ví dụ~1 đến ví dụ~8), cả hai
chỉ số cùng giảm --- Jaccard từ $0.83$ về $0.00$ và cosine TF-IDF từ $0.88$ về $0.00$ ---
và khả năng phát hiện của TF-IDF suy giảm tương ứng.

\begin{table}[htbp]
\centering
\scriptsize
\renewcommand{\arraystretch}{1.3}
\begin{tabular}{@{}p{0.35cm}p{1.9cm}p{3.9cm}p{3.9cm}p{0.8cm}p{0.9cm}p{1.15cm}@{}}
\toprule
\textbf{\#} & \textbf{Kỹ thuật} & \textbf{Câu nguồn} & \textbf{Câu đạo văn} & \textbf{Jaccard} & \textbf{cos TF-IDF} & \textbf{Phát hiện} \\
\midrule
1 & Sao chép nguyên văn & \emph{\dots velocity fluctuations dominate the small-scale structure observed in spectral-line data cubes.} & \emph{\dots velocity fluctuations dominate the small-scale structure seen in spectral-line data cubes.} & 0.83 & 0.88 & Có (mạnh) \\
2 & Thay từ đồng nghĩa (nhẹ) & \emph{Deep neural networks achieve state-of-the-art accuracy on large-scale image classification benchmarks.} & \emph{Deep neural networks attain leading performance on large-scale image recognition benchmarks.} & 0.44 & 0.54 & Có \\
3 & Thay từ đồng nghĩa (mạnh) & \emph{The proposed algorithm significantly reduces the computational cost of training large language models.} & \emph{Our method markedly lowers the training expense of building sizeable language models.} & 0.19 & 0.21 & Yếu \\
4 & Đảo trật tự câu & \emph{Because the dataset is imbalanced, we apply oversampling before training the classifier.} & \emph{We apply oversampling before training the classifier, because the dataset is imbalanced.} & 1.00 & 1.00 & Có (không phụ thuộc thứ tự) \\
5 & Chắp vá (mosaic) & \emph{Transformer architectures rely on self-attention to model long-range dependencies in sequences.} & \emph{To capture long-range dependencies, sequence models built on transformers use the self-attention mechanism.} & 0.28 & 0.40 & Yếu \\
6 & Dịch đi--dịch lại & \emph{The experimental results confirm that our approach outperforms existing baselines by a wide margin.} & \emph{Findings from the experiments verify a substantial performance gap in favor of the presented technique over prior methods.} & 0.00 & 0.00 & Không \\
7 & Đạo ý & \emph{Increasing the number of training epochs improves accuracy but eventually leads to overfitting.} & \emph{Prolonged optimization raises predictive quality up to a point, after which the model starts memorizing the data.} & 0.00 & 0.00 & Không \\
8 & Đạo ý (cực) & \emph{Regularization techniques such as dropout prevent co-adaptation of neurons and improve generalization.} & \emph{By randomly disabling units during learning, one discourages units from relying on each other and yields models that transfer better to unseen inputs.} & 0.00 & 0.00 & Không \\
\bottomrule
\end{tabular}
\caption{Suy giảm độ trùng từ vựng qua các kỹ thuật đạo văn, minh hoạ ở mức câu. Jaccard được
tính trên tập từ nội dung (chữ thường, loại stop-word và từ $\le 2$ ký tự); \emph{cos TF-IDF}
là cosine giữa hai vector TF-IDF theo đúng cấu hình tầng truy hồi (\texttt{sublinear\_tf},
loại stop-word tiếng Anh).}
\label{tab:paraphrase}
\end{table}


\section{Truy hồi nguồn (Source Retrieval)}

\subsection{Bài toán và phương pháp}

Với mỗi tài liệu nghi vấn, việc đối chiếu chi tiết trực tiếp với toàn bộ $7\,949$
tài liệu nguồn là bất khả thi về chi phí tính toán. Tầng truy hồi có nhiệm vụ thu hẹp
kho nguồn về một tập nhỏ các ứng viên khả nghi: đầu vào là tài liệu nghi vấn cùng chỉ mục
của kho, đầu ra là danh sách $k$ nguồn được xếp hạng theo mức tương đồng.

Phương pháp được sử dụng là TF-IDF kết hợp độ đo cosine. Mỗi tài liệu được biểu diễn bằng
một vector TF-IDF trên $20\,000$ ký tự đầu (cắt theo ký tự, mỗi tài liệu một vector), với
cấu hình \texttt{max\_features} $=100\,000$, \texttt{sublinear\_tf}, và loại bỏ stop-word
tiếng Anh. Vector của tài liệu nghi vấn được so cosine với toàn bộ ma trận nguồn để lấy
$k$ tài liệu có điểm cao nhất.

\subsection{Ví dụ minh hoạ biểu diễn TF-IDF}

Để làm rõ cách TF-IDF đánh trọng số cho từ, xét một kho đồ chơi gồm bốn văn bản ngắn:

\begin{itemize}
  \item[$d_1$] (truy vấn): \emph{deep neural networks classify images with high accuracy}
  \item[$d_2$:] \emph{neural networks are trained on large datasets}
  \item[$d_3$:] \emph{the model reduces the computational cost of training}
  \item[$d_4$:] \emph{convolutional networks improve image classification accuracy}
\end{itemize}

Sau khi loại stop-word tiếng Anh, mỗi từ được gán trọng số theo công thức của thư viện
\texttt{scikit-learn}:
\[
\mathrm{idf}(t) = \ln\frac{1+N}{1+\mathrm{df}(t)} + 1,
\qquad
w(t,d) = \big(1 + \ln \mathrm{tf}(t,d)\big)\cdot \mathrm{idf}(t),
\]
trong đó $N=4$ là số văn bản, $\mathrm{df}(t)$ là số văn bản chứa từ $t$, và $\mathrm{tf}(t,d)$ là
tần suất từ $t$ trong văn bản $d$; vector cuối được chuẩn hoá $L_2$. Bảng~\ref{tab:tfidf-demo}
liệt kê trọng số TF-IDF của câu truy vấn $d_1$.

\begin{table}[htbp]
\centering
\renewcommand{\arraystretch}{1.2}
\begin{tabular}{lccc}
\toprule
\textbf{Từ} & \textbf{df} & \textbf{idf} & \textbf{TF-IDF ($d_1$)} \\
\midrule
\emph{classify} & 1 & 1.916 & \textbf{0.421} \\
\emph{deep}     & 1 & 1.916 & \textbf{0.421} \\
\emph{images}   & 1 & 1.916 & \textbf{0.421} \\
\emph{high}     & 1 & 1.916 & \textbf{0.421} \\
\emph{accuracy} & 2 & 1.511 & 0.332 \\
\emph{neural}   & 2 & 1.511 & 0.332 \\
\emph{networks} & 3 & 1.223 & 0.269 \\
\bottomrule
\end{tabular}
\caption{Trọng số TF-IDF của câu truy vấn $d_1$ trên kho đồ chơi bốn văn bản ($N=4$),
sắp theo trọng số giảm dần.}
\label{tab:tfidf-demo}
\end{table}

Vì mỗi từ chỉ xuất hiện một lần trong $d_1$, thành phần $\mathrm{tf}$ là đồng đều, nên chênh
lệch trọng số hoàn toàn do IDF quyết định. Các từ đặc trưng chỉ có mặt ở một văn bản
(\emph{classify}, \emph{deep}, \emph{images}, \emph{high}; $\mathrm{df}=1$) nhận trọng số cao
nhất ($0.421$), trong khi từ phổ biến \emph{networks} ($\mathrm{df}=3$) chỉ còn $0.269$. Nhờ
cơ chế này, độ đo cosine giữa hai tài liệu bị chi phối bởi các \emph{thuật ngữ đặc trưng dùng
chung}; đó chính là tín hiệu giúp tầng truy hồi xác định đúng nguồn.

\subsection{Lựa chọn TF-IDF và kết quả so sánh}

TF-IDF được lựa chọn sau khi so sánh trực tiếp với phương pháp biểu diễn ngữ nghĩa (neural
embedding) trên tập validation gồm $5\,522$ tài liệu nghi vấn. Kết quả được trình bày trong
Bảng~\ref{tab:retrieval-recall}: TF-IDF vượt trội ở mọi mức Recall@k cũng như MRR.

\begin{table}[htbp]
\centering
\renewcommand{\arraystretch}{1.2}
\begin{tabular}{lcccc}
\toprule
\textbf{Phương pháp} & \textbf{R@1} & \textbf{R@5} & \textbf{R@10} & \textbf{MRR} \\
\midrule
\textbf{TF-IDF (từ vựng)}       & \textbf{0.975} & \textbf{0.994} & \textbf{0.995} & \textbf{0.985} \\
Neural (bge-small)              & 0.908          & 0.965          & 0.975          & 0.934          \\
\bottomrule
\end{tabular}
\caption{Recall@k và MRR trên tập validation ($5\,522$ tài liệu nghi vấn).}
\label{tab:retrieval-recall}
\end{table}

Sự vượt trội của TF-IDF xuất phát từ hai yếu tố. Thứ nhất, đạo văn trong bộ PAN có độ trùng
lặp từ vựng cao, và các tài liệu khoa học chia sẻ nhiều thuật ngữ chuyên ngành, tên riêng,
ký hiệu và con số; đây là các tín hiệu mà mô hình túi từ (bag-of-words) nắm bắt trực tiếp.
Thứ hai, biểu diễn ngữ nghĩa \texttt{bge-small} bị chi phối bởi chủ đề: các tài liệu cùng
lĩnh vực nằm gần nhau trong không gian embedding, làm mờ ranh giới giữa nguồn thật và các
tài liệu chỉ trùng chủ đề, dẫn đến hạng của nguồn đúng bị hạ thấp.

\subsection{Giới hạn của biểu diễn từ vựng}

Do là mô hình túi từ, TF-IDF chỉ nắm bắt phần từ vựng bề mặt được dùng chung giữa hai văn
bản. Cần phân biệt hai dạng che giấu (obfuscation) với ảnh hưởng khác nhau:

\begin{itemize}
  \item \textbf{Đảo trật tự câu hoặc mệnh đề} không ảnh hưởng đến TF-IDF, vì biểu diễn túi từ
  bất biến với thứ tự từ: toàn bộ từ được giữ nguyên nên độ trùng từ vựng không đổi (ví dụ~4
  trong Bảng~\ref{tab:paraphrase}, Jaccard $=1.00$).
  \item \textbf{Thay từ đồng nghĩa, dịch đi--dịch lại (back-translation) và viết lại ý} là giới
  hạn thực sự: khi các từ dùng chung bị thay bằng từ khác về mặt hình thức, độ trùng từ vựng
  giảm dần về $0$ và tín hiệu lexical biến mất. Ở các đoạn bị che giấu theo cách này, TF-IDF
  không còn khả năng phát hiện.
\end{itemize}

Bảng~\ref{tab:paraphrase} minh hoạ mức suy giảm độ trùng từ vựng qua các kỹ thuật đạo văn,
đo bằng hai chỉ số: Jaccard trên tập từ nội dung, và cosine giữa hai vector TF-IDF (cùng cấu
hình với tầng truy hồi). Khi từ vựng dùng chung bị thay thế dần (ví dụ~1 đến ví dụ~8), cả hai
chỉ số cùng giảm --- Jaccard từ $0.83$ về $0.00$ và cosine TF-IDF từ $0.88$ về $0.00$ ---
và khả năng phát hiện của TF-IDF suy giảm tương ứng.

\begin{table}[htbp]
\centering
\scriptsize
\renewcommand{\arraystretch}{1.3}
\begin{tabular}{@{}p{0.35cm}p{1.9cm}p{3.9cm}p{3.9cm}p{0.8cm}p{0.9cm}p{1.15cm}@{}}
\toprule
\textbf{\#} & \textbf{Kỹ thuật} & \textbf{Câu nguồn} & \textbf{Câu đạo văn} & \textbf{Jaccard} & \textbf{cos TF-IDF} & \textbf{Phát hiện} \\
\midrule
1 & Sao chép nguyên văn & \emph{\dots velocity fluctuations dominate the small-scale structure observed in spectral-line data cubes.} & \emph{\dots velocity fluctuations dominate the small-scale structure seen in spectral-line data cubes.} & 0.83 & 0.88 & Có (mạnh) \\
2 & Thay từ đồng nghĩa (nhẹ) & \emph{Deep neural networks achieve state-of-the-art accuracy on large-scale image classification benchmarks.} & \emph{Deep neural networks attain leading performance on large-scale image recognition benchmarks.} & 0.44 & 0.54 & Có \\
3 & Thay từ đồng nghĩa (mạnh) & \emph{The proposed algorithm significantly reduces the computational cost of training large language models.} & \emph{Our method markedly lowers the training expense of building sizeable language models.} & 0.19 & 0.21 & Yếu \\
4 & Đảo trật tự câu & \emph{Because the dataset is imbalanced, we apply oversampling before training the classifier.} & \emph{We apply oversampling before training the classifier, because the dataset is imbalanced.} & 1.00 & 1.00 & Có (không phụ thuộc thứ tự) \\
5 & Chắp vá (mosaic) & \emph{Transformer architectures rely on self-attention to model long-range dependencies in sequences.} & \emph{To capture long-range dependencies, sequence models built on transformers use the self-attention mechanism.} & 0.28 & 0.40 & Yếu \\
6 & Dịch đi--dịch lại & \emph{The experimental results confirm that our approach outperforms existing baselines by a wide margin.} & \emph{Findings from the experiments verify a substantial performance gap in favor of the presented technique over prior methods.} & 0.00 & 0.00 & Không \\
7 & Đạo ý & \emph{Increasing the number of training epochs improves accuracy but eventually leads to overfitting.} & \emph{Prolonged optimization raises predictive quality up to a point, after which the model starts memorizing the data.} & 0.00 & 0.00 & Không \\
8 & Đạo ý (cực) & \emph{Regularization techniques such as dropout prevent co-adaptation of neurons and improve generalization.} & \emph{By randomly disabling units during learning, one discourages units from relying on each other and yields models that transfer better to unseen inputs.} & 0.00 & 0.00 & Không \\
\bottomrule
\end{tabular}
\caption{Suy giảm độ trùng từ vựng qua các kỹ thuật đạo văn, minh hoạ ở mức câu. Jaccard được
tính trên tập từ nội dung (chữ thường, loại stop-word và từ $\le 2$ ký tự); \emph{cos TF-IDF}
là cosine giữa hai vector TF-IDF theo đúng cấu hình tầng truy hồi (\texttt{sublinear\_tf},
loại stop-word tiếng Anh).}
\label{tab:paraphrase}
\end{table}


\section{Truy hồi nguồn (Source Retrieval)}

\subsection{Bài toán và phương pháp}

Với mỗi tài liệu nghi vấn, việc đối chiếu chi tiết trực tiếp với toàn bộ $7\,949$
tài liệu nguồn là bất khả thi về chi phí tính toán. Tầng truy hồi có nhiệm vụ thu hẹp
kho nguồn về một tập nhỏ các ứng viên khả nghi: đầu vào là tài liệu nghi vấn cùng chỉ mục
của kho, đầu ra là danh sách $k$ nguồn được xếp hạng theo mức tương đồng.

Phương pháp được sử dụng là TF-IDF kết hợp độ đo cosine. Mỗi tài liệu được biểu diễn bằng
một vector TF-IDF trên $20\,000$ ký tự đầu (cắt theo ký tự, mỗi tài liệu một vector), với
cấu hình \texttt{max\_features} $=100\,000$, \texttt{sublinear\_tf}, và loại bỏ stop-word
tiếng Anh. Vector của tài liệu nghi vấn được so cosine với toàn bộ ma trận nguồn để lấy
$k$ tài liệu có điểm cao nhất.

\subsection{Ví dụ minh hoạ biểu diễn TF-IDF}

Để làm rõ cách TF-IDF đánh trọng số cho từ, xét một kho đồ chơi gồm bốn văn bản ngắn:

\begin{itemize}
  \item[$d_1$] (truy vấn): \emph{deep neural networks classify images with high accuracy}
  \item[$d_2$:] \emph{neural networks are trained on large datasets}
  \item[$d_3$:] \emph{the model reduces the computational cost of training}
  \item[$d_4$:] \emph{convolutional networks improve image classification accuracy}
\end{itemize}

Sau khi loại stop-word tiếng Anh, mỗi từ được gán trọng số theo công thức của thư viện
\texttt{scikit-learn}:
\[
\mathrm{idf}(t) = \ln\frac{1+N}{1+\mathrm{df}(t)} + 1,
\qquad
w(t,d) = \big(1 + \ln \mathrm{tf}(t,d)\big)\cdot \mathrm{idf}(t),
\]
trong đó $N=4$ là số văn bản, $\mathrm{df}(t)$ là số văn bản chứa từ $t$, và $\mathrm{tf}(t,d)$ là
tần suất từ $t$ trong văn bản $d$; vector cuối được chuẩn hoá $L_2$. Bảng~\ref{tab:tfidf-demo}
liệt kê trọng số TF-IDF của câu truy vấn $d_1$.

\begin{table}[htbp]
\centering
\renewcommand{\arraystretch}{1.2}
\begin{tabular}{lccc}
\toprule
\textbf{Từ} & \textbf{df} & \textbf{idf} & \textbf{TF-IDF ($d_1$)} \\
\midrule
\emph{classify} & 1 & 1.916 & \textbf{0.421} \\
\emph{deep}     & 1 & 1.916 & \textbf{0.421} \\
\emph{images}   & 1 & 1.916 & \textbf{0.421} \\
\emph{high}     & 1 & 1.916 & \textbf{0.421} \\
\emph{accuracy} & 2 & 1.511 & 0.332 \\
\emph{neural}   & 2 & 1.511 & 0.332 \\
\emph{networks} & 3 & 1.223 & 0.269 \\
\bottomrule
\end{tabular}
\caption{Trọng số TF-IDF của câu truy vấn $d_1$ trên kho đồ chơi bốn văn bản ($N=4$),
sắp theo trọng số giảm dần.}
\label{tab:tfidf-demo}
\end{table}

Vì mỗi từ chỉ xuất hiện một lần trong $d_1$, thành phần $\mathrm{tf}$ là đồng đều, nên chênh
lệch trọng số hoàn toàn do IDF quyết định. Các từ đặc trưng chỉ có mặt ở một văn bản
(\emph{classify}, \emph{deep}, \emph{images}, \emph{high}; $\mathrm{df}=1$) nhận trọng số cao
nhất ($0.421$), trong khi từ phổ biến \emph{networks} ($\mathrm{df}=3$) chỉ còn $0.269$. Nhờ
cơ chế này, độ đo cosine giữa hai tài liệu bị chi phối bởi các \emph{thuật ngữ đặc trưng dùng
chung}; đó chính là tín hiệu giúp tầng truy hồi xác định đúng nguồn.

\subsection{Lựa chọn TF-IDF và kết quả so sánh}

TF-IDF được lựa chọn sau khi so sánh trực tiếp với phương pháp biểu diễn ngữ nghĩa (neural
embedding) trên tập validation gồm $5\,522$ tài liệu nghi vấn. Kết quả được trình bày trong
Bảng~\ref{tab:retrieval-recall}: TF-IDF vượt trội ở mọi mức Recall@k cũng như MRR.

\begin{table}[htbp]
\centering
\renewcommand{\arraystretch}{1.2}
\begin{tabular}{lcccc}
\toprule
\textbf{Phương pháp} & \textbf{R@1} & \textbf{R@5} & \textbf{R@10} & \textbf{MRR} \\
\midrule
\textbf{TF-IDF (từ vựng)}       & \textbf{0.975} & \textbf{0.994} & \textbf{0.995} & \textbf{0.985} \\
Neural (bge-small)              & 0.908          & 0.965          & 0.975          & 0.934          \\
\bottomrule
\end{tabular}
\caption{Recall@k và MRR trên tập validation ($5\,522$ tài liệu nghi vấn).}
\label{tab:retrieval-recall}
\end{table}

Sự vượt trội của TF-IDF xuất phát từ hai yếu tố. Thứ nhất, đạo văn trong bộ PAN có độ trùng
lặp từ vựng cao, và các tài liệu khoa học chia sẻ nhiều thuật ngữ chuyên ngành, tên riêng,
ký hiệu và con số; đây là các tín hiệu mà mô hình túi từ (bag-of-words) nắm bắt trực tiếp.
Thứ hai, biểu diễn ngữ nghĩa \texttt{bge-small} bị chi phối bởi chủ đề: các tài liệu cùng
lĩnh vực nằm gần nhau trong không gian embedding, làm mờ ranh giới giữa nguồn thật và các
tài liệu chỉ trùng chủ đề, dẫn đến hạng của nguồn đúng bị hạ thấp.

\subsection{Giới hạn của biểu diễn từ vựng}

Do là mô hình túi từ, TF-IDF chỉ nắm bắt phần từ vựng bề mặt được dùng chung giữa hai văn
bản. Cần phân biệt hai dạng che giấu (obfuscation) với ảnh hưởng khác nhau:

\begin{itemize}
  \item \textbf{Đảo trật tự câu hoặc mệnh đề} không ảnh hưởng đến TF-IDF, vì biểu diễn túi từ
  bất biến với thứ tự từ: toàn bộ từ được giữ nguyên nên độ trùng từ vựng không đổi (ví dụ~4
  trong Bảng~\ref{tab:paraphrase}, Jaccard $=1.00$).
  \item \textbf{Thay từ đồng nghĩa, dịch đi--dịch lại (back-translation) và viết lại ý} là giới
  hạn thực sự: khi các từ dùng chung bị thay bằng từ khác về mặt hình thức, độ trùng từ vựng
  giảm dần về $0$ và tín hiệu lexical biến mất. Ở các đoạn bị che giấu theo cách này, TF-IDF
  không còn khả năng phát hiện.
\end{itemize}

Bảng~\ref{tab:paraphrase} minh hoạ mức suy giảm độ trùng từ vựng qua các kỹ thuật đạo văn,
đo bằng hai chỉ số: Jaccard trên tập từ nội dung, và cosine giữa hai vector TF-IDF (cùng cấu
hình với tầng truy hồi). Khi từ vựng dùng chung bị thay thế dần (ví dụ~1 đến ví dụ~8), cả hai
chỉ số cùng giảm --- Jaccard từ $0.83$ về $0.00$ và cosine TF-IDF từ $0.88$ về $0.00$ ---
và khả năng phát hiện của TF-IDF suy giảm tương ứng.

\begin{table}[htbp]
\centering
\scriptsize
\renewcommand{\arraystretch}{1.3}
\begin{tabular}{@{}p{0.35cm}p{1.9cm}p{3.9cm}p{3.9cm}p{0.8cm}p{0.9cm}p{1.15cm}@{}}
\toprule
\textbf{\#} & \textbf{Kỹ thuật} & \textbf{Câu nguồn} & \textbf{Câu đạo văn} & \textbf{Jaccard} & \textbf{cos TF-IDF} & \textbf{Phát hiện} \\
\midrule
1 & Sao chép nguyên văn & \emph{\dots velocity fluctuations dominate the small-scale structure observed in spectral-line data cubes.} & \emph{\dots velocity fluctuations dominate the small-scale structure seen in spectral-line data cubes.} & 0.83 & 0.88 & Có (mạnh) \\
2 & Thay từ đồng nghĩa (nhẹ) & \emph{Deep neural networks achieve state-of-the-art accuracy on large-scale image classification benchmarks.} & \emph{Deep neural networks attain leading performance on large-scale image recognition benchmarks.} & 0.44 & 0.54 & Có \\
3 & Thay từ đồng nghĩa (mạnh) & \emph{The proposed algorithm significantly reduces the computational cost of training large language models.} & \emph{Our method markedly lowers the training expense of building sizeable language models.} & 0.19 & 0.21 & Yếu \\
4 & Đảo trật tự câu & \emph{Because the dataset is imbalanced, we apply oversampling before training the classifier.} & \emph{We apply oversampling before training the classifier, because the dataset is imbalanced.} & 1.00 & 1.00 & Có (không phụ thuộc thứ tự) \\
5 & Chắp vá (mosaic) & \emph{Transformer architectures rely on self-attention to model long-range dependencies in sequences.} & \emph{To capture long-range dependencies, sequence models built on transformers use the self-attention mechanism.} & 0.28 & 0.40 & Yếu \\
6 & Dịch đi--dịch lại & \emph{The experimental results confirm that our approach outperforms existing baselines by a wide margin.} & \emph{Findings from the experiments verify a substantial performance gap in favor of the presented technique over prior methods.} & 0.00 & 0.00 & Không \\
7 & Đạo ý & \emph{Increasing the number of training epochs improves accuracy but eventually leads to overfitting.} & \emph{Prolonged optimization raises predictive quality up to a point, after which the model starts memorizing the data.} & 0.00 & 0.00 & Không \\
8 & Đạo ý (cực) & \emph{Regularization techniques such as dropout prevent co-adaptation of neurons and improve generalization.} & \emph{By randomly disabling units during learning, one discourages units from relying on each other and yields models that transfer better to unseen inputs.} & 0.00 & 0.00 & Không \\
\bottomrule
\end{tabular}
\caption{Suy giảm độ trùng từ vựng qua các kỹ thuật đạo văn, minh hoạ ở mức câu. Jaccard được
tính trên tập từ nội dung (chữ thường, loại stop-word và từ $\le 2$ ký tự); \emph{cos TF-IDF}
là cosine giữa hai vector TF-IDF theo đúng cấu hình tầng truy hồi (\texttt{sublinear\_tf},
loại stop-word tiếng Anh).}
\label{tab:paraphrase}
\end{table}


\section{Truy hồi nguồn (Source Retrieval)}

\subsection{Bài toán và phương pháp}

Với mỗi tài liệu nghi vấn, việc đối chiếu chi tiết trực tiếp với toàn bộ $7\,949$
tài liệu nguồn là bất khả thi về chi phí tính toán. Tầng truy hồi có nhiệm vụ thu hẹp
kho nguồn về một tập nhỏ các ứng viên khả nghi: đầu vào là tài liệu nghi vấn cùng chỉ mục
của kho, đầu ra là danh sách $k$ nguồn được xếp hạng theo mức tương đồng.

Phương pháp được sử dụng là TF-IDF kết hợp độ đo cosine. Mỗi tài liệu được biểu diễn bằng
một vector TF-IDF trên $20\,000$ ký tự đầu (cắt theo ký tự, mỗi tài liệu một vector), với
cấu hình \texttt{max\_features} $=100\,000$, \texttt{sublinear\_tf}, và loại bỏ stop-word
tiếng Anh. Vector của tài liệu nghi vấn được so cosine với toàn bộ ma trận nguồn để lấy
$k$ tài liệu có điểm cao nhất.

\subsection{Ví dụ minh hoạ biểu diễn TF-IDF}

Để làm rõ cách TF-IDF đánh trọng số cho từ, xét một kho đồ chơi gồm bốn văn bản ngắn:

\begin{itemize}
  \item[$d_1$] (truy vấn): \emph{deep neural networks classify images with high accuracy}
  \item[$d_2$:] \emph{neural networks are trained on large datasets}
  \item[$d_3$:] \emph{the model reduces the computational cost of training}
  \item[$d_4$:] \emph{convolutional networks improve image classification accuracy}
\end{itemize}

Sau khi loại stop-word tiếng Anh, mỗi từ được gán trọng số theo công thức của thư viện
\texttt{scikit-learn}:
\[
\mathrm{idf}(t) = \ln\frac{1+N}{1+\mathrm{df}(t)} + 1,
\qquad
w(t,d) = \big(1 + \ln \mathrm{tf}(t,d)\big)\cdot \mathrm{idf}(t),
\]
trong đó $N=4$ là số văn bản, $\mathrm{df}(t)$ là số văn bản chứa từ $t$, và $\mathrm{tf}(t,d)$ là
tần suất từ $t$ trong văn bản $d$; vector cuối được chuẩn hoá $L_2$. Bảng~\ref{tab:tfidf-demo}
liệt kê trọng số TF-IDF của câu truy vấn $d_1$.

\begin{table}[htbp]
\centering
\renewcommand{\arraystretch}{1.2}
\begin{tabular}{lccc}
\toprule
\textbf{Từ} & \textbf{df} & \textbf{idf} & \textbf{TF-IDF ($d_1$)} \\
\midrule
\emph{classify} & 1 & 1.916 & \textbf{0.421} \\
\emph{deep}     & 1 & 1.916 & \textbf{0.421} \\
\emph{images}   & 1 & 1.916 & \textbf{0.421} \\
\emph{high}     & 1 & 1.916 & \textbf{0.421} \\
\emph{accuracy} & 2 & 1.511 & 0.332 \\
\emph{neural}   & 2 & 1.511 & 0.332 \\
\emph{networks} & 3 & 1.223 & 0.269 \\
\bottomrule
\end{tabular}
\caption{Trọng số TF-IDF của câu truy vấn $d_1$ trên kho đồ chơi bốn văn bản ($N=4$),
sắp theo trọng số giảm dần.}
\label{tab:tfidf-demo}
\end{table}

Vì mỗi từ chỉ xuất hiện một lần trong $d_1$, thành phần $\mathrm{tf}$ là đồng đều, nên chênh
lệch trọng số hoàn toàn do IDF quyết định. Các từ đặc trưng chỉ có mặt ở một văn bản
(\emph{classify}, \emph{deep}, \emph{images}, \emph{high}; $\mathrm{df}=1$) nhận trọng số cao
nhất ($0.421$), trong khi từ phổ biến \emph{networks} ($\mathrm{df}=3$) chỉ còn $0.269$. Nhờ
cơ chế này, độ đo cosine giữa hai tài liệu bị chi phối bởi các \emph{thuật ngữ đặc trưng dùng
chung}; đó chính là tín hiệu giúp tầng truy hồi xác định đúng nguồn.

\subsection{Lựa chọn TF-IDF và kết quả so sánh}

TF-IDF được lựa chọn sau khi so sánh trực tiếp với phương pháp biểu diễn ngữ nghĩa (neural
embedding) trên tập validation gồm $5\,522$ tài liệu nghi vấn. Kết quả được trình bày trong
Bảng~\ref{tab:retrieval-recall}: TF-IDF vượt trội ở mọi mức Recall@k cũng như MRR.

\begin{table}[htbp]
\centering
\renewcommand{\arraystretch}{1.2}
\begin{tabular}{lcccc}
\toprule
\textbf{Phương pháp} & \textbf{R@1} & \textbf{R@5} & \textbf{R@10} & \textbf{MRR} \\
\midrule
\textbf{TF-IDF (từ vựng)}       & \textbf{0.975} & \textbf{0.994} & \textbf{0.995} & \textbf{0.985} \\
Neural (bge-small)              & 0.908          & 0.965          & 0.975          & 0.934          \\
\bottomrule
\end{tabular}
\caption{Recall@k và MRR trên tập validation ($5\,522$ tài liệu nghi vấn).}
\label{tab:retrieval-recall}
\end{table}

Sự vượt trội của TF-IDF xuất phát từ hai yếu tố. Thứ nhất, đạo văn trong bộ PAN có độ trùng
lặp từ vựng cao, và các tài liệu khoa học chia sẻ nhiều thuật ngữ chuyên ngành, tên riêng,
ký hiệu và con số; đây là các tín hiệu mà mô hình túi từ (bag-of-words) nắm bắt trực tiếp.
Thứ hai, biểu diễn ngữ nghĩa \texttt{bge-small} bị chi phối bởi chủ đề: các tài liệu cùng
lĩnh vực nằm gần nhau trong không gian embedding, làm mờ ranh giới giữa nguồn thật và các
tài liệu chỉ trùng chủ đề, dẫn đến hạng của nguồn đúng bị hạ thấp.

\subsection{Giới hạn của biểu diễn từ vựng}

Do là mô hình túi từ, TF-IDF chỉ nắm bắt phần từ vựng bề mặt được dùng chung giữa hai văn
bản. Cần phân biệt hai dạng che giấu (obfuscation) với ảnh hưởng khác nhau:

\begin{itemize}
  \item \textbf{Đảo trật tự câu hoặc mệnh đề} không ảnh hưởng đến TF-IDF, vì biểu diễn túi từ
  bất biến với thứ tự từ: toàn bộ từ được giữ nguyên nên độ trùng từ vựng không đổi (ví dụ~4
  trong Bảng~\ref{tab:paraphrase}, Jaccard $=1.00$).
  \item \textbf{Thay từ đồng nghĩa, dịch đi--dịch lại (back-translation) và viết lại ý} là giới
  hạn thực sự: khi các từ dùng chung bị thay bằng từ khác về mặt hình thức, độ trùng từ vựng
  giảm dần về $0$ và tín hiệu lexical biến mất. Ở các đoạn bị che giấu theo cách này, TF-IDF
  không còn khả năng phát hiện.
\end{itemize}

Bảng~\ref{tab:paraphrase} minh hoạ mức suy giảm độ trùng từ vựng qua các kỹ thuật đạo văn,
đo bằng hai chỉ số: Jaccard trên tập từ nội dung, và cosine giữa hai vector TF-IDF (cùng cấu
hình với tầng truy hồi). Khi từ vựng dùng chung bị thay thế dần (ví dụ~1 đến ví dụ~8), cả hai
chỉ số cùng giảm --- Jaccard từ $0.83$ về $0.00$ và cosine TF-IDF từ $0.88$ về $0.00$ ---
và khả năng phát hiện của TF-IDF suy giảm tương ứng.

\begin{table}[htbp]
\centering
\scriptsize
\renewcommand{\arraystretch}{1.3}
\begin{tabular}{@{}p{0.35cm}p{1.9cm}p{3.9cm}p{3.9cm}p{0.8cm}p{0.9cm}p{1.15cm}@{}}
\toprule
\textbf{\#} & \textbf{Kỹ thuật} & \textbf{Câu nguồn} & \textbf{Câu đạo văn} & \textbf{Jaccard} & \textbf{cos TF-IDF} & \textbf{Phát hiện} \\
\midrule
1 & Sao chép nguyên văn & \emph{\dots velocity fluctuations dominate the small-scale structure observed in spectral-line data cubes.} & \emph{\dots velocity fluctuations dominate the small-scale structure seen in spectral-line data cubes.} & 0.83 & 0.88 & Có (mạnh) \\
2 & Thay từ đồng nghĩa (nhẹ) & \emph{Deep neural networks achieve state-of-the-art accuracy on large-scale image classification benchmarks.} & \emph{Deep neural networks attain leading performance on large-scale image recognition benchmarks.} & 0.44 & 0.54 & Có \\
3 & Thay từ đồng nghĩa (mạnh) & \emph{The proposed algorithm significantly reduces the computational cost of training large language models.} & \emph{Our method markedly lowers the training expense of building sizeable language models.} & 0.19 & 0.21 & Yếu \\
4 & Đảo trật tự câu & \emph{Because the dataset is imbalanced, we apply oversampling before training the classifier.} & \emph{We apply oversampling before training the classifier, because the dataset is imbalanced.} & 1.00 & 1.00 & Có (không phụ thuộc thứ tự) \\
5 & Chắp vá (mosaic) & \emph{Transformer architectures rely on self-attention to model long-range dependencies in sequences.} & \emph{To capture long-range dependencies, sequence models built on transformers use the self-attention mechanism.} & 0.28 & 0.40 & Yếu \\
6 & Dịch đi--dịch lại & \emph{The experimental results confirm that our approach outperforms existing baselines by a wide margin.} & \emph{Findings from the experiments verify a substantial performance gap in favor of the presented technique over prior methods.} & 0.00 & 0.00 & Không \\
7 & Đạo ý & \emph{Increasing the number of training epochs improves accuracy but eventually leads to overfitting.} & \emph{Prolonged optimization raises predictive quality up to a point, after which the model starts memorizing the data.} & 0.00 & 0.00 & Không \\
8 & Đạo ý (cực) & \emph{Regularization techniques such as dropout prevent co-adaptation of neurons and improve generalization.} & \emph{By randomly disabling units during learning, one discourages units from relying on each other and yields models that transfer better to unseen inputs.} & 0.00 & 0.00 & Không \\
\bottomrule
\end{tabular}
\caption{Suy giảm độ trùng từ vựng qua các kỹ thuật đạo văn, minh hoạ ở mức câu. Jaccard được
tính trên tập từ nội dung (chữ thường, loại stop-word và từ $\le 2$ ký tự); \emph{cos TF-IDF}
là cosine giữa hai vector TF-IDF theo đúng cấu hình tầng truy hồi (\texttt{sublinear\_tf},
loại stop-word tiếng Anh).}
\label{tab:paraphrase}
\end{table}
